\chapter{Introduction}

\section{What is \amlet?}

\amlet{} is an acronym for \underline{A}nother \underline{M}ixed
\underline{L}ogit \underline{E}stimation \underline{T}ool.
As its name suggests, it is a software mainly designed to estimate various
kind of mixed logit models.
It can be used either as a standalone program, either as a C library.
\amlet{} is an open source project (see Appendix~\ref{sec:license};
the last version can be downloaded at the address
\url{http://amlet.slashbin.net}.

\section{Supported systems}

\amlet{} development is primarily Unix-oriented, while the software
 can also be used on Windows systems, with some technical
 limitations\footnote{In particular, the timings will be less
 accurate.}.
Most of the development and tests have been performed on Linux systems.
The software should also runs without major trouble on BSD systems, or Mac OS X.
Windows users should be able to compile AMLET using MinGW, available at \url{http://www.mingw.org}, at the expense of some loss of performance.
For Windows users, it is recommended to install Linux in a virtual machine, using for instance VirtualBox, available for free at \url{http://www.virtualbox.org}. See Section~\ref{sec:Windows} for more details.

Unless otherwise stated, we will assume in this document that you
work on a Unix-like system; the explications are still valid for the
other environment. We also assume that you use \verb|bash| as your
shell, since it is the default choice for major Linux distributions,
as well as \verb|mingw|; some users could have to slighty adapt the
explanation if they use csh shells, and its variants.
If you do not know what is your current shell, you probably already
use \verb|bash|.
It can also be determined by reading the file \verb|/etc/password|.

\section{Requirements}

\subsection{Hardware requirements}

There is no formal minimum requirements for \amlet{}, since they
crucially depend on the size of dataset size and the number of random
variables that are present in the model.
We however recommand a minimum of 1Go RAM, knowing that the available memory constrains the achievable accuracy.
\amlet{} should work smoothly on any modern computer.

\subsection{Software requirements}

Various software elements are required in order to correctly install
\amlet{}, as well as some optional elements, if you want to
install additional tools along with the main package.

\begin{description}
\item[C compiler]
We recommand to use the compiler delivered with your
computer; on Linux-machine and mingw, this is typically the GNU C
Compiler (\gcc).
\item[Compilers]
The code is mainly written in C and C++, but contains portions written in Fortran, mainly imported from LAPACK.
Consequently, both a C and a Fortran compilers are requires; the code has been successfully tested with \verb|gcc| and \verb|gfortran|, available in any Linux distribution.
\item[Unix shell]
In order to easily perform the installation, autotools
(\textsf{autoconf}, \textsf{automake}, \textsf{libtoolize}) are highly
recommanded.
The detailed installation process can however often be pursued
without them.
If you are using Mac OS X, you can experience problems with
\verb|libtoolize|. Using the development environment proposed in the
third-party software collections fink
(\url{http://www.finkproject.org/}) can resolve this issue.
\item[Lex, Yacc]
Driver parsing with \amlet{} requires standard lexical analyser and parser.
\amlet{} requires \textsf{flex}, and \textsf{yacc} or \textsf{bison}.
\item[Perl]
Some optional analysis tools provided with \amlet{} rely on Perl scripting
language.
\item[Oratio, Ophelia]
These are two companion libraries, that have to be installed first and can be downloaded on \url{http://amlet.slashbin.net}.
\end{description}

\subsection{Windows users}
\label{sec:Windows}

The most convenient solution is to install Linux on a Virtual Machine, using for instance VirtualBox (\url{http://www.virtualbox.org}).
A virtual machine allows to emulate a PC, allowing the installation of a complete operating system.
There are of course some overheads in executation time, as it means that one operating system is running inside another one, but a Virtual Machine can be easily uninstalled if no more needed, without affecting the host Operating System.
It is therefore the easiest solution.

A more complex solution is to install \mingw{}, standing for Minimal Gnu system for Windows.
\mingw{} can be download at the address \url{http://www.mingw.org}.
Pay attention to install all the development packages during the installation, in order to have to required tools to compile the code.
It is also a good idea to install a correct text editor. Various
editors are available under Windows: \textsf{Vim}, \textsf{Emacs},
\textsf{Xemacs}.

\section{Acknowledgements}

Like for many projects, the development of \amlet{} has only been
possible with the help of many people.
First of all, I would like to express my gratitude so Cinzia Cirillo,
for the hard work that we have achieved together. This software could
not has been written without her advices, based on her huge experience
of the practical aspects of the problem. I would also like to thank
Philippe L. Toint, who was my Ph.D. advisor during four years, and
teached me the most important aspects of nonlinear
programming. Finally, my thanks go to Belgian National Fund for
Scientific Research, who gave me the grant that made the initial work
on this project possible.
The development of \amlet{} has been
initiated at the University of Namur (Belgium), continued in the
Imperial College London (United Kingdom) and Cerfacs (Toulouse, France).
I pursue it during my (very limited) free time at the University of Montreal (Canada).

The documentation of a software is always a difficult exercise. I
would therefore like to thank anyone that helped us to improve the
present document. In particular, I would like to express my
gratitude to Roberta ``Speedy'' Pellican\`o for her suggestions,
during my study stay in the University Frederico II of Naples, Italy
(June 2004), and Jasper Knockaert.
The software has also benefit from many suggestions during its
development. The following list is certainly far from be complete; my
apologies go to each one I have forgotten. I will cite here (by
alphabetical order) Michel Bierlaire and St\'e{}phane Hess.

Finally, the painting reproduced the front cover is due to Pieter van
Steenwyck, under the name {\sl Ars longa, vitta brevis}.

\section{Installation}
\label{sec:installation}

\amlet{} is currently distributed in source code only, so that the
installation process primarily involves the compilation of the various
source files.
\amlet{} first requires the installation of the libraries \oratio{} and \ophelia{}. We refer to their respective documentations for details about the corresponding installation procedures.

As for the libraries, you can compile \amlet{} with the commands
\begin{verbatim}
./configure
make
sudo make install
\end{verbatim}
Optimization compilation options can be specified as arguments of \verb|configure|, for instance:
\begin{verbatim}
./configure CFLAGS=''-O2'' FFLAGS="-O2"
\end{verbatim}

